\documentclass[11pt]{amsproc}

\usepackage[a4paper,margin=2.6cm]{geometry}
\usepackage{amsmath,amssymb,amsthm}
\usepackage[T1]{fontenc}
\usepackage[utf8]{inputenc}

\newtheorem*{problema}{Problema}

\newcommand{\R}{\mathbb{R}}
\newcommand{\Z}{\mathbb{Z}}
\newcommand{\Q}{\mathbb{Q}}
\newcommand{\ip}[2]{\langle #1,\, #2\rangle}

\begin{document}

\section*{Problema}

Una casa de apuestas ofrece $n$ apuestas distintas sobre una carrera con $r$
resultados posibles. Apostar un peso a la apuesta $j$ devuelve $m_{ij}$ pesos si
ocurre el resultado $i$, donde todos los $m_{ij}$ son números enteros; un retorno
negativo significa una pérdida de esa cantidad. Se permite arriesgar un monto entero
$c_j$ en cada apuesta $j$, y un monto negativo significa tomar el lado de la casa en
esa apuesta. Así, si ocurre el resultado $i$, la ganancia neta es
\[
  \sum_{j=1}^{n} m_{ij}\,c_j .
\]

Diremos que un vector de pesos reales $(b_1,\dots,b_r)\neq(0,\dots,0)$ es
\emph{equilibrante} si, al ponderar los resultados con $b_1,\dots,b_r$, toda apuesta
queda pareja:
\[
  \sum_{i=1}^{r} b_i\, m_{ij}=0
  \qquad\text{para todo } j=1,\dots,n .
\]

\begin{problema}
Supongamos que en todo vector equilibrante hay al menos dos pesos de signos
opuestos. Probar que existen montos enteros $c_1,\dots,c_n$ que garantizan una
ganancia estrictamente positiva, cualquiera sea el resultado de la carrera.
\end{problema}

\section*{Solución}

Escribamos $M=(m_{ij})$ para la matriz $r\times n$ de retornos, de modo que un vector
de montos $c\in\R^n$ produce el vector de ganancias $Mc\in\R^r$. Sea
\[
  V=\{\,Mc\ :\ c\in\R^n\,\}\subseteq\R^r
\]
el espacio columna de $M$. Como $b^{\mathsf T}M=0$ vale exactamente cuando $b$ es
ortogonal a $Mc$ para todo $c$, los vectores equilibrantes son precisamente los
elementos no nulos de $V^{\perp}$.

La hipótesis se reformula entonces así:
\begin{equation}\label{eq:hyp}
  \text{si } b\in V^{\perp},\ b\neq 0, \text{ entonces } b \text{ tiene alguna coordenada negativa.}
\end{equation}
En efecto, un $b\in V^\perp$ no nulo con todas sus coordenadas $\ge 0$ sería un
vector equilibrante sin dos pesos de signos opuestos.

Hay que producir $c\in\Z^n$ con todas las coordenadas de $Mc$ estrictamente
positivas. Primero encontramos un tal $c$ en $\R^n$, y después arreglamos la
integralidad sin costo alguno.

\subsection*{Paso 1}

Sean $e_1,\dots,e_r$ la base canónica de $\R^r$ y sea $P\colon\R^r\to\R^r$ la
proyección ortogonal sobre $V$; recordemos que $P=P^{\mathsf T}=P^2$ y
$\ker P=V^{\perp}$. Sea $\Delta=\{\,b\in\R^r:\ b_i\ge 0,\ \sum_i b_i=1\,\}$ el
símplex estándar y definamos
\[
  A \;=\; P(\Delta) \;=\; \Bigl\{\ \sum_{i=1}^{r} b_i\,Pe_i \ :\ b_i\ge 0,\ \sum_i b_i=1 \ \Bigr\}
  \;\subseteq\; V .
\]
Por ser imagen del conjunto compacto y convexo $\Delta$ por una aplicación lineal,
$A$ es compacto y convexo. Además
\[
  0\notin A,
\]
porque $Pb=0$ para algún $b\in\Delta$ significaría $b\in V^{\perp}$ con $b\ge 0$ y
$b\neq 0$, contradiciendo \eqref{eq:hyp}.

Por compacidad, la función continua $u\mapsto\|u\|$ alcanza un mínimo sobre $A$,
digamos en $v\in A$; y $v\neq 0$. Afirmamos que
\begin{equation}\label{eq:var}
  \ip{v}{u}\ \ge\ \|v\|^{2}\qquad\text{para todo } u\in A .
\end{equation}
En efecto, fijemos $u\in A$. Por convexidad $v+t(u-v)\in A$ para todo $t\in(0,1]$, de
modo que, por la minimalidad de $\|v\|$,
\[
  \|v\|^{2}\ \le\ \|v+t(u-v)\|^{2}
  \ =\ \|v\|^{2}+2t\,\ip{v}{u-v}+t^{2}\|u-v\|^{2}.
\]
Cancelando $\|v\|^2$, dividiendo por $t>0$ y haciendo $t\to 0^{+}$ se obtiene
$\ip{v}{u-v}\ge 0$, que es \eqref{eq:var}.

Apliquemos ahora \eqref{eq:var} al punto $u=Pe_i\in A$. Como $v\in V$ se tiene
$Pv=v$, y como $P$ es autoadjunta,
\[
  \ip{v}{Pe_i}=\ip{Pv}{e_i}=\ip{v}{e_i}=v_i .
\]
Por lo tanto
\[
  v_i\ \ge\ \|v\|^{2}\ >\ 0\qquad\text{para todo } i=1,\dots,r,
\]
así que el vector $v\in V$ tiene \emph{todas} sus coordenadas estrictamente
positivas. Eligiendo $c^{\ast}\in\R^n$ con $Mc^{\ast}=v$, esos montos (reales) ya
garantizan una ganancia de al menos $\|v\|^2$ en cada resultado.

\subsection*{Paso 2} 

Consideremos
\[
  U=\{\,c\in\R^{n}\ :\ (Mc)_i>0 \text{ para todo } i=1,\dots,r\,\}.
\]
La aplicación $c\mapsto Mc$ es continua y el ortante positivo abierto de $\R^r$ es
abierto, así que $U$ es abierto; por el Paso 1 es no vacío. Como $\Q^{n}$ es denso en
$\R^{n}$, el abierto $U$ contiene un punto racional $q$. Sea $N$ un denominador común
positivo de $q_1,\dots,q_n$ y pongamos $c=Nq$. Entonces $c\in\Z^{n}$ y
\[
  Mc = N\,(Mq),
\]
cuyas coordenadas son todas estrictamente positivas. Estos montos enteros ganan,
pase lo que pase. 

\end{document}
